% Generated draft: requires mathematical and source review.
\documentclass[11pt]{article}
\usepackage[margin=25mm]{geometry}
\usepackage{fontspec}
\setmainfont{Latin Modern Roman}
\usepackage{amsmath,amssymb,mathtools,unicode-math}
\setmathfont{Latin Modern Math}
\usepackage{enumitem,needspace,xurl,hyperref}
\setlength{\parindent}{0pt}
\setlength{\parskip}{5pt}
\setlength{\emergencystretch}{3em}
\newcommand{\N}{\mathbb{N}}
\newcommand{\Z}{\mathbb{Z}}
\newcommand{\Q}{\mathbb{Q}}
\newcommand{\R}{\mathbb{R}}
\newcommand{\C}{\mathbb{C}}
\newcommand{\F}{\mathbb{F}}
\newcommand{\A}{\mathbb{A}}
\newcommand{\PP}{\mathbb{P}}
\newcommand{\E}{\mathbb{E}}
\newcommand{\eps}{\varepsilon}
\newcommand{\dd}{\,\mathrm{d}}
\begin{document}
\section*{UnsolvedMath ID: 30004326}
\textbf{Source problem code:} catalogue entry
\subsection{Definitions and mathematical statement}
Let \(\pi:X\to\mathbb{P}^{2}_{k}\) be the blowup at \(r\) distinct very general points \(p_1,\ldots,p_r\). Write \(H\) for the pullback of the class of a line and \(E_i\) for the exceptional divisor above \(p_i\). The source does not specify the base field \(k\) or its characteristic. The usual geometric meaning of ``very general'' is outside a countable union of proper closed subsets of the parameter space of ordered distinct point configurations; the field hypotheses needed for this interpretation remain unspecified.

Let
\[
D=dH-\sum_{i=1}^{r}m_iE_i.
\]
For an integral divisor \(D\), its coefficients \(d,m_i\) are integers. The source does not explicitly impose nonnegativity or give their intended ranges. Let \(|D|\) denote the complete linear system of effective divisors linearly equivalent to \(D\), with
\[
\dim|D|=\dim_k H^0(X,\mathcal{O}_X(D))-1,
\]
where \(\dim\varnothing=-1\). Define the displayed expected dimension by
\[
e(D)=\max\left\{\binom{d+2}{2}-\sum_{i=1}^{r}\binom{m_i+1}{2}-1,-1\right\}.
\]
The binomial terms have their ordinary combinatorial meaning when \(d,m_i\geq0\); their intended interpretation outside that range is not specified by the source. Call \(|D|\) special, in the sense of this record, when \(\dim|D|>e(D)\).

For an effective divisor \(A=\sum_j a_jC_j\), with distinct prime divisors \(C_j\) and positive integer coefficients \(a_j\), nonreduced means that some \(a_j\geq2\). The supplied question is whether, for every choice of data in the intended scope,
\[
\dim|D|>e(D)
\quad\Longrightarrow\quad
\forall A\in|D|\;\exists j\;(a_j\geq2).
\]
Here the prime components and coefficients refer to each individual divisor \(A\) on \(X\). The extracted statement gives no restriction on \(r\), whereas the background describes the SHGH conjecture for at least ten points. Thus whether the intended question assumes \(r\geq10\), together with the ground-field and coefficient hypotheses, requires clarification before this becomes a fully specified conjecture.
\subsection{Short English statement}
The question asks whether every effective divisor in a special complete linear system on a plane blowup at very general points must contain a repeated irreducible component. The source leaves the ground field and numerical scope incomplete, and its background mentions \(r\geq10\) although the statement does not.
\subsection{Sources}
\begin{itemize}
\item Local UnsolvedMath registry, permanent ID \(30004326\), title \emph{Nonreduced Special Linear Systems on Plane Blowups}, source problem code \texttt{catalogue entry}, collection \emph{Oberwolfach Reports Open Problems}. No individual proposer is supplied; the registry gives proposed year 2019.
\item Supplied report citation: \emph{Mini-Workshop: Seshadri Constants} (2020), \emph{Oberwolfach Reports}, pp.\ 3271--3313, DOI \texttt{10.4171/owr/2019/53}. \url{https://doi.org/10.4171/owr/2019/53}. The background describes these reports as documenting mathematical workshops and their open-problem discussions.
\item Supplied literature reference: Roe, \emph{On the Nagata conjecture}. \url{https://arxiv.org/abs/math/0304124}.
\item Supplied literature reference: \emph{Rationality of Seshadri constants on general blow-ups}. \url{https://www.sciencedirect.com/science/article/pii/S0022404920300426}.
\item Supplied literature reference: \emph{Oberwolfach Reports, SHGH conjecture}. \url{https://doi.org/10.4171/owr/2019/53}.
\item Literature assessment dated exactly \texttt{2026-08-22}: the supplied curation labels the problem open and describes the extracted sentence as introducing the SHGH conjecture for at least ten points, whose general assertion remains unresolved and underlies open Seshadri problems. This is supplied dated curation, not an independently verified literature assessment; the record explicitly requests independent verification of time-sensitive status and cited claims.
\item The registry labels the statement \texttt{corrected\_verified}, but explains that only formula delimiters, notation, grammar, and typography were repaired, with mathematical content and scope preserved. This label does not establish mathematical correctness or source fidelity beyond the supplied qualification.
\end{itemize}
\end{document}
